\section{Related Work}

The landscape of audio signal restoration has traditionally relied on supervised learning paradigms requiring paired clean-noisy data \cite{snoek2012}. However, recent advances in unsupervised reconstruction have shifted the focus toward leveraging corrupted observations directly. Lehtinen et al.\ introduced Noise2Noise\cite{lehtinen2018}, demonstrating that statistical reasoning allows models to learn restoration mappings without access to clean ground truth by exploiting noise correlations across samples. This principle is critical for our target domain, where periodic seam-wrap artifacts in wavetable synthesis preclude the availability of pristine reference signals \cite{stoller2018}. Concurrently, population-based training methods have emerged as a robust alternative to hyperparameter sensitivity issues inherent in standard deep learning pipelines\cite{jaderberg2017}. By maintaining an ensemble of networks and updating weights via averaging rather than backpropagation through the entire dataset, Population Based Training (PBT) facilitates rapid adaptation to diverse acoustic environments. We integrate this mechanism into our outer loop optimization strategy to stabilize training dynamics during the exploration phase \cite{schulman2017}.

In the realm of automated architecture discovery, evolutionary algorithms have proven effective for navigating vast design spaces where gradient-based search is computationally prohibitive\cite{elsken2019}. Real et al.\ proposed Regularized Evolution (RealNas)\cite{real2019}, which applies genetic operators to neural network topologies while regularizing the objective function against overfitting. This approach aligns with our hybrid reinforcement learning and evolutionary strategy, where a policy gradient controller guides an agent through architectural modifications \cite{khadka2018}. Furthermore, Pham et al.\ demonstrated that parameter sharing across candidate subgraphs significantly reduces computational overhead during architecture search\cite{pham2018}, enabling the evaluation of thousands of configurations within feasible timeframes. These techniques are essential for constructing a flexible backbone capable of handling non-stationary noise distributions in real-time audio processing tasks \cite{elsken2019}.

Differentiable Digital Signal Processing (DDSP) offers another avenue by embedding domain-specific priors directly into the neural network architecture\cite{engel2020}. Unlike generic black-box models, DDSP utilizes explicit representations of sound generation processes such as harmonic oscillators and noise sources. While our current work focuses on denoising periodic artifacts rather than full synthesis, we adopt similar principles to constrain the solution space toward physically plausible waveforms \cite{stoller2018}. The integration of Mixture-of-Experts layers\cite{shazeer2017} further supports this goal by allowing distinct sub-networks to specialize in different frequency bands or artifact types. Although we do not employ Generative Adversarial Networks (GANs) for speech enhancement \cite{pascual2017}, the adversarial training concepts inform our loss function design, ensuring that reconstructed signals adhere strictly to spectral constraints imposed by the underlying wavetable physics.

Despite these advancements, several critical gaps remain in existing literature regarding periodic artifact removal and meta-learning strategies tailored to audio synthesis. Behl et al.\ explored Model-Agnostic Meta-Learning (MAML) for fast adaptation\cite{finn2017}, yet their approach assumes independent task distributions that do not capture the continuous nature of wavetable parameter spaces \cite{jaderberg2017}. Similarly, while Differentiable Architecture Search methods like DARTS have achieved state-of-the-art results in image classification\cite{liu2019}, they often neglect the computational efficiency required for real-time audio applications. In the specific domain of source separation and denoising, recent works such as Demucs \cite{dfossez2019} rely heavily on magnitude spectrograms which discard phase information crucial for reconstructing periodic waveforms accurately\cite{stoller2018}. Diffusion models like DiffWave offer versatile synthesis capabilities but lack the interpretability needed to isolate specific artifact types without extensive fine-tuning \cite{kong2021}. Additionally, hybrid approaches combining spectral and waveform domains have shown promise in music separation tasks \cite{dfossez2021}, yet they typically assume stationary noise characteristics that fail under dynamic playback conditions.

Our proposed methodology distinguishes itself by explicitly addressing the challenge of periodic seam-wrap crackle through a unified framework combining reinforcement learning, evolutionary algorithms, population-based training, and neural architecture search\cite{jaderberg2017,khadka2018,real2019,schulman2017}. Unlike prior efforts that treat denoising as a static regression problem or rely on fixed architectures optimized offline \cite{elsken2019}, our system dynamically adapts its topology and hyperparameters during inference to match the specific characteristics of incoming audio streams. This adaptability is particularly important when dealing with varying playback speeds or pitch shifts common in wavetable synthesis, where traditional methods often introduce audible artifacts due to misalignment between learned filters and actual signal content \cite{engel2020}. By leveraging unsupervised learning principles akin to Noise2Noise\cite{lehtinen2018}, we ensure robust performance even when clean reference data is unavailable. Moreover, the incorporation of Mixture-of-Experts layers allows our model to selectively engage specialized sub-networks depending on whether the input contains tonal components or broadband noise \cite{shazeer2017}.

Current experiments are ongoing with an extended run targeting 250,000 iterations using a hybrid PPO+GA+PBT+NAS configuration incorporating depth scaling and MoE structures. Preliminary results indicate early champions achieving residual approximations near 0.987 against historical baselines around 0.7 \cite{stoller2018}. These findings suggest significant potential for improving upon existing benchmarks while maintaining computational efficiency suitable for deployment in resource-constrained environments such as portable synthesizers or embedded audio processors \cite{jaderberg2017}. Future iterations will focus on refining the reward shaping mechanisms to better penalize phase distortions introduced by aggressive filtering operations.
